\section{Methodology}

\subsection{Problem Setting}

We evaluate neural decoders for the planted clique problem on random graphs $G(n,1/2,k)$. Each instance is generated by first sampling an Erd\H{o}s--R\'{e}nyi graph $G(n,1/2)$, where every possible edge is included independently with probability $1/2$, and then planting a clique $C^\star$ of size $k$. The goal is to recover a large valid clique, ideally the planted clique. We report three metrics:
\begin{itemize}
  \item \textbf{Approximation}: $|C_{\mathrm{found}}|/|C^\star|$.
  \item \textbf{Overlap}: $|C_{\mathrm{found}}\cap C^\star|/|C^\star|$.
  \item \textbf{Exact}: the fraction of instances for which the recovered clique equals the planted clique.
\end{itemize}

All reported experiments use $n=1000$ and follow the planted-clique evaluation protocol of the Learning-to-Rank setup. We divide instances into three regimes according to the planted clique size $k$: \textbf{Easy} uses $k\in[62,100]$, \textbf{Medium} uses $k\in[36,61]$, and \textbf{Hard} uses $k\in[20,35]$. Easy instances often admit a strong static ranking signal, Medium instances typically require iterative residual pruning, and Hard instances contain a weaker planted signal with many dense distractor subgraphs.

\subsection{Model Variants}

We compare two GNN architectures: a standard message-passing GNN (SGNN) and a model-aware GNN (MAGNN). Both models start from the same constant node feature,
$
x_v^{(0)}=1,
$
so any useful ranking signal must be inferred from the graph structure, except for the explicit objective-gradient channel available only to MAGNN.

\paragraph{Standard GNN (SGNN).}
The input is embedded as
\[
  h_v^{(0)} = W_{\mathrm{in}}x_v^{(0)} + b_{\mathrm{in}} .
\]
At internal step $\ell$, each vertex aggregates messages from its graph neighbors,
\[
  q_v^{(\ell)}
  =
  \operatorname{Agg}_{u\in N(v)} h_u^{(\ell)} .
\]
We use normalized sum-style aggregation rather than pure mean aggregation. With identical initial features, mean aggregation can remove differences in neighborhood size, whereas sum-style aggregation preserves such structural variation while keeping the scale controlled.

The SGNN update is
\[
  h_v^{(\ell+1)}
  =
  \sigma\!\left(
    M_{\mathrm{upd}}\!\left(h_v^{(\ell)},q_v^{(\ell)}\right)
  \right),
\]
where $M_{\mathrm{upd}}$ is a shared multilayer perceptron. After $L$ internal steps, the model outputs a logit
\[
  z_v = W_{\mathrm{out}}h_v^{(L)} + b_{\mathrm{out}},
\]
which is used directly as the keep-score:
\[
  s_v=z_v .
\]
Higher scores indicate vertices that should be kept.

\paragraph{Model-Aware GNN (MAGNN).}
MAGNN uses the same backbone, aggregation rule, nonlinearity, and output logit definition as SGNN. The only architectural difference is that MAGNN injects a per-node objective-gradient channel during message passing.

At internal step $\ell$, MAGNN forms an intermediate logit
\[
  z_v^{(\ell)}
  =
  W_{\mathrm{out}}h_v^{(\ell)} + b_{\mathrm{out}},
\]
maps it to a relaxed selection variable
\[
  p_v^{(\ell)}
  =
  \operatorname{sigmoid}\!\left(z_v^{(\ell)}\right),
\]
and computes
\[
  g_v^{(\ell)}
  =
  \frac{\partial \mathcal{L}_{\mathrm{clique}}(p^{(\ell)};A,k)}
  {\partial p_v^{(\ell)}} .
\]
The hidden-state update becomes
\[
  h_v^{(\ell+1)}
  =
  \sigma\!\left(
    M_{\mathrm{upd}}\!\left(h_v^{(\ell)},q_v^{(\ell)},g_v^{(\ell)}\right)
  \right).
\]
The sigmoid only maps unconstrained logits to the domain of the continuous relaxation; the model-aware signal is the gradient $g_v^{(\ell)}$. Thus, SGNN and MAGNN differ only in whether the current objective gradient is exposed as an internal per-node signal.

In all experiments, both models use hidden dimension $64$, $L=10$ shared internal message-passing steps, and $\sigma=\tanh$.

\subsection{Clique Objective and Gradient}

Let $z\in\mathbb{R}^n$ denote the model logits and let
$
p=\operatorname{sigmoid}(z)\in[0,1]^n
$
be the relaxed vertex-selection vector. Let $A$ be the adjacency matrix and let $B$ be the non-edge adjacency matrix, with zero diagonal:
\[
  B_{ij}=1-A_{ij}, \qquad B_{ii}=0 .
\]
We use the differentiable clique objective
\[
  \mathcal{L}_{\mathrm{clique}}(p;A,k)
  =
  -\frac{p^\top A p}{k(k-1)}
  +\frac{p^\top B p}{k(n-k)} .
\]
The first term rewards placing mass on edges, while the second penalizes placing mass on non-edges.

The gradient injected by MAGNN is
\[
  \nabla_p \mathcal{L}_{\mathrm{clique}}(p;A,k)
  =
  -\frac{2Ap}{k(k-1)}
  +\frac{2Bp}{k(n-k)} .
\]
This exposes a local marginal signal of the current relaxed objective.

\subsection{Training Regimes}

We use two training regimes.

\paragraph{Objective Training.}
The model is trained only with the unsupervised clique objective:
\[
  \mathcal{L}_{\mathrm{train}}
  =
  \mathcal{L}_{\mathrm{clique}} .
\]
This tests whether the architecture and objective alone induce useful vertex scores.

\paragraph{Self-Predictor LDR Training.}
The self-predictor regime trains a cheap incremental decoder to imitate an adaptive LDR-style pruning process. In LDR pruning, the algorithm repeatedly scores the current residual graph, removes the least promising vertex, and continues until the remaining vertices form a clique. This adaptive process is strong because scores are recomputed after each pruning step, but it is expensive when the scoring function is a GNN.

We use this adaptive process as a teacher. At step $t$, the teacher reruns the GNN on the current residual graph $G_t$ and removes
\[
  v_t^{\mathrm{teacher}}
  =
  \arg\min_{v\in V(G_t)} s_v(G_t).
\]
The student runs the GNN only once on the original graph and then follows the teacher's removal sequence using incrementally updated scores $\tilde{s}^{(t)}$. The self-imitation loss is
\[
  \mathcal{L}_{\mathrm{self}}
  =
  \frac{1}{T}\sum_{t=0}^{T-1}
  \operatorname{CE}\!\left(
    \operatorname{softmax}(-\tilde{s}^{(t)}),
    v_t^{\mathrm{teacher}}
  \right),
\]
where the minus sign reflects that lower scores correspond to removed vertices. The total loss is
\[
  \mathcal{L}_{\mathrm{train}}
  =
  \mathcal{L}_{\mathrm{clique}}
  +
  \lambda_{\mathrm{self}}\mathcal{L}_{\mathrm{self}} .
\]
In the reported experiments, $\lambda_{\mathrm{self}}=1$.

\paragraph{Training Pseudocode.}
\begin{algorithmic}[1]
  \State Sample a planted-clique instance $(G,k)$. \# get one training graph
  \State $\tilde{s} \leftarrow \mathrm{GNN}(G)$. \# run the student once
  \State $G_R \leftarrow G$. \# initialize residual graph
  \State $\mathcal{L}_{\mathrm{self}}\leftarrow 0$. \# initialize self-imitation loss

  \For{$t=1,\ldots,T$}
  \State $s^T \leftarrow \mathrm{GNN}(G_R)$. \# teacher reruns GNN on current residual graph
  \State $v \leftarrow \arg\min_{u\in V(G_R)} s^T_u$. \# teacher removes lowest-scoring vertex
  \State $\mathcal{L}_{\mathrm{self}} \leftarrow \mathcal{L}_{\mathrm{self}} + \operatorname{CE}\!\left(\operatorname{softmax}(-\tilde{s}_{V(G_R)}),v\right)$. \# train student to predict teacher removal
  \State $G_R \leftarrow G_R \setminus \{v\}$. \# update residual graph
  \State $\tilde{s}_i \leftarrow \tilde{s}_i + U_\theta\!\left(\tilde{s}_i,\tilde{s}_v,A_{iv},g_i,g_v\right)$ for all $i\in V(G_R)$. \# learned score update
  \EndFor

  \State $\mathcal{L}_{\mathrm{self}} \leftarrow \frac{1}{T}\mathcal{L}_{\mathrm{self}}$. \# average over pruning steps
  \State $\mathcal{L}_{\mathrm{train}} = \mathcal{L}_{\mathrm{clique}} + \lambda_{\mathrm{self}}\mathcal{L}_{\mathrm{self}}$. \# combine losses
  \State Optimize $\mathcal{L}_{\mathrm{train}}$. \# update parameters
\end{algorithmic}

\subsection{Learned Incremental Updater}

During incremental decoding, the GNN is run once to produce initial vertex scores. After each removed vertex $v$, the remaining scores are updated by a learned nodewise updater:
\[
  \Delta s_i
  =
  U_\theta\!\left(s_i,s_v,A_{iv},g_i,g_v\right),
  \qquad
  s_i \leftarrow s_i + \Delta s_i ,
  \qquad i\in V(G_R).
\]
Here $G_R$ is the current residual graph, $s_i$ is the current score of vertex $i$, $s_v$ is the score of the removed vertex, and $A_{iv}$ indicates whether $i$ was adjacent to the removed vertex. For SGNN, the gradient entries $g_i,g_v$ are set to zero. For MAGNN, they are objective-gradient values.

The updater is applied independently to each remaining vertex. Therefore, each pruning step costs $O(n)$, and the full incremental decoder consists of one $O(n^2)$ GNN pass followed by $O(n)$ learned score updates. This yields an $O(n^2)$ residual decoder while still allowing scores to adapt after each pruning decision.

\subsection{Decoding Regimes}

We evaluate two decoders in the main experiments.

\paragraph{Rerun-Pruned Decoder.}
The rerun-pruned decoder repeatedly reruns the GNN on the current residual graph, removes the lowest-scoring vertex, stops once the remaining vertices form a clique, and then tries to add removed vertices back whenever doing so preserves clique feasibility. This is the expensive adaptive baseline.

\begin{algorithmic}[1]
  \State $C \leftarrow V(G)$. \# current candidate vertex set
  \State $H \leftarrow \emptyset$. \# set of removed vertices, stored in removal order

  \While{$C$ does not form a clique}
  \State $s \leftarrow \mathrm{GNN}$ applied to the subgraph induced by $C$. \# score current candidates
  \State $v \leftarrow \arg\min_{u\in C} s_u$. \# select least promising candidate
  \State $C \leftarrow C\setminus\{v\}$. \# prune selected vertex
  \State Add $v$ to the end of $H$. \# remember removal order
  \EndWhile

  \For{$v$ from last to first in $H$}
  \If{$C\cup\{v\}$ forms a clique}
  \State $C \leftarrow C\cup\{v\}$. \# restore feasible vertex
  \EndIf
  \EndFor

  \State \Return $C$. \# recovered clique
\end{algorithmic}

\paragraph{Incremental Decoder.}
The incremental decoder runs the GNN only once. It then repeatedly removes the lowest-scoring vertex and updates all remaining scores using the learned updater. It uses the same stopping and add-back rule as rerun-pruned decoding.

\begin{algorithmic}[1]
  \State $s \leftarrow \mathrm{GNN}(G)$. \# compute initial scores once
  \State $C \leftarrow V(G)$. \# current candidate vertex set
  \State $H \leftarrow \emptyset$. \# removed vertices, stored in removal order

  \While{$C$ does not form a clique}
  \State $v \leftarrow \arg\min_{u\in C} s_u$. \# select least promising candidate
  \State $C \leftarrow C\setminus\{v\}$. \# prune selected vertex
  \State Add $v$ to the end of $H$. \# remember removal order
  \State $s_i \leftarrow s_i + U_\theta(s_i,s_v,A_{iv},g_i,g_v)$ for all $i\in C$. \# update remaining scores
  \EndWhile

  \For{$v$ from last to first in $H$}
  \If{$C\cup\{v\}$ forms a clique}
  \State $C \leftarrow C\cup\{v\}$. \# restore feasible vertex
  \EndIf
  \EndFor

  \State \Return $C$. \# recovered clique
\end{algorithmic}

On dense graphs, a single GNN pass costs $O(n^2)$. Rerun-pruned decoding may require $O(n)$ such passes, yielding $O(n^3)$ time. Incremental decoding instead uses one $O(n^2)$ GNN pass followed by $O(n)$ removals, each with an $O(n)$ learned update, giving $O(n^2)$ total decoding time.

\subsection{Representation Diagnostics}

We additionally analyze the final node embeddings by projecting them onto their first principal component. These plots are diagnostic rather than part of the main decoder. They test whether the learned embedding geometry organizes vertices along a degree- or rank-like direction.

\begin{figure}[t]
  \centering
  \fbox{\parbox{0.92\linewidth}{\centering Placeholder: SGNN embedding PCA colored by degree/rank.}}
  \caption{Placeholder for SGNN embedding geometry. Nodes will be plotted using the first two principal components of their final embeddings and colored by degree or degree rank.}
  \label{fig:sgnn-pca-placeholder}
\end{figure}

\begin{figure}[t]
  \centering
  \fbox{\parbox{0.92\linewidth}{\centering Placeholder: MAGNN embedding PCA colored by degree/rank.}}
  \caption{Placeholder for MAGNN embedding geometry. Nodes will be plotted using the first two principal components of their final embeddings and colored by degree or degree rank.}
  \label{fig:magnn-pca-placeholder}
\end{figure}

\section{Results}

\subsection{Main Performance Results}

Table~\ref{tab:main-results} reports planted clique recovery for the four main model variants: SGNN and MAGNN, each trained either with objective training or self-predictor LPR training. The table reports the rerun-pruned decoder and the learned incremental decoder.

\begin{table*}[t]
  \centering
  \small
  \begin{tabular}{ll|ccc|ccc|ccc}
    \toprule
    & & \multicolumn{3}{c|}{Easy} & \multicolumn{3}{c|}{Medium} & \multicolumn{3}{c}{Hard} \\
    Model & Decoder & Approx & Overlap & Exact & Approx & Overlap & Exact & Approx & Overlap & Exact \\
    \midrule
    \multirow{2}{*}{SGNN + Objective}
    & Rerun-Pruned & 1.000 & 1.000 & 1.000 & 1.000 & 1.000 & 1.000 & 0.621 & 0.342 & 0.200 \\
    & Incremental & 0.354 & 0.282 & 0.200 & 0.285 & 0.132 & 0.000 & 0.390 & 0.023 & 0.000 \\
    \midrule
    \multirow{2}{*}{MAGNN + Objective}
    & Rerun-Pruned & 1.000 & 1.000 & 1.000 & 1.000 & 1.000 & 1.000 & 0.607 & 0.339 & 0.200 \\
    & Incremental & 0.928 & 0.925 & 0.800 & 0.772 & 0.740 & 0.600 & 0.370 & 0.114 & 0.000 \\
    \midrule
    \multirow{2}{*}{SGNN + Self-Predictor}
    & Rerun-Pruned & 1.000 & 1.000 & 1.000 & 1.000 & 1.000 & 1.000 & 0.806 & 0.698 & 0.600 \\
    & Incremental & 0.116 & 0.003 & 0.000 & 0.195 & 0.004 & 0.000 & 0.373 & 0.008 & 0.000 \\
    \midrule
    \multirow{2}{*}{MAGNN + Self-Predictor}
    & Rerun-Pruned & 1.000 & 1.000 & 1.000 & 0.698 & 0.600 & 0.600 & 0.390 & 0.031 & 0.000 \\
    & Incremental & 1.000 & 1.000 & 1.000 & 1.000 & 1.000 & 1.000 & 0.689 & 0.427 & 0.400 \\
    \bottomrule
  \end{tabular}
  \caption{Main planted clique recovery results on $n=1000$. Approx is $|C_{\mathrm{found}}|/|C^\star|$, overlap is $|C_{\mathrm{found}}\cap C^\star|/|C^\star|$, and exact is the exact planted-clique recovery rate.}
  \label{tab:main-results}
\end{table*}

\subsection{Interpretation}

The results reproduce the learning-to-rank phenomenon and extend it with a model-aware incremental decoder. First, the rerun-pruned decoder is strong for SGNN, reaching perfect recovery on Easy and Medium instances under objective training. This mirrors the standard learning-to-rank behavior: a GNN can learn a useful vertex ranking, and repeated pruning on residual graphs can convert this ranking into high-quality clique recovery.

However, the SGNN ranking does not transfer to cheap incremental decoding. Under objective training, SGNN incremental decoding reaches only $0.354$ approximation on Easy and $0.285$ on Medium. Under self-predictor training, SGNN incremental decoding remains weak: $0.116$ on Easy and $0.195$ on Medium. This shows that good rerun-pruned behavior does not imply that the learned scores can be updated cheaply after each removal.

MAGNN changes this behavior. With objective training alone, MAGNN incremental decoding reaches $0.928$ approximation on Easy and $0.772$ on Medium, compared with SGNN's $0.354$ and $0.285$ under the same training regime. With self-predictor training, MAGNN incremental decoding reaches perfect recovery on Easy and Medium and substantially improves Hard recovery to $0.689$ approximation. This supports the central claim: the objective-gradient channel provides local optimization information that a learned updater can use to approximate rerun-style residual decoding without rerunning the GNN.

The Hard regime remains substantially more difficult. Even the best MAGNN incremental model reaches $0.689$ approximation and $0.400$ exact recovery. This is expected: in the Hard regime, the planted signal is weak and many dense distractor subgraphs exist. Nonetheless, the MAGNN incremental decoder is the strongest incremental method in the table.

\subsection{One-Pass Diagnostic Results}

Although the main decoder comparison uses rerun-pruned and incremental decoding, one-pass results are useful diagnostics. They show whether the GNN learns a useful static ranking before any residual adaptation.

\begin{table*}[t]
  \centering
  \small
  \begin{tabular}{l|ccc|ccc|ccc}
    \toprule
    & \multicolumn{3}{c|}{Easy} & \multicolumn{3}{c|}{Medium} & \multicolumn{3}{c}{Hard} \\
    Model & Approx & Overlap & Exact & Approx & Overlap & Exact & Approx & Overlap & Exact \\
    \midrule
    SGNN + Objective & 0.964 & 0.963 & 0.900 & 0.692 & 0.637 & 0.400 & 0.422 & 0.046 & 0.000 \\
    MAGNN + Objective & 1.000 & 1.000 & 1.000 & 0.518 & 0.455 & 0.200 & 0.485 & 0.151 & 0.100 \\
    SGNN + Self-Predictor & 1.000 & 1.000 & 1.000 & 0.505 & 0.438 & 0.200 & 0.375 & 0.118 & 0.000 \\
    MAGNN + Self-Predictor & 0.963 & 0.961 & 0.900 & 0.810 & 0.797 & 0.600 & 0.448 & 0.140 & 0.000 \\
    \bottomrule
  \end{tabular}
  \caption{One-pass diagnostic performance. These results show the quality of the static ranking produced by each trained model before rerunning or incremental updates.}
  \label{tab:one-pass-diagnostics}
\end{table*}

The one-pass results confirm that all models learn meaningful static rankings on Easy instances. Medium and Hard instances require more adaptive behavior. The best Medium one-pass result among the main models is MAGNN with self-predictor training, reaching $0.810$ approximation and $0.600$ exact recovery. However, the main gain of MAGNN is not only static ranking: it is the ability to support learned incremental residual updates.

\subsection{PCA and Degree-Rank Diagnostics}

We compute PCA on the final node embeddings after one full-graph GNN pass and report the variance explained by the first two principal components and Spearman correlation between original graph degree and PC1. The provided logs report Spearman correlations; Pearson correlations were not included in the current run output and should be added in the next run if needed.

\begin{table*}[t]
  \centering
  \small
  \begin{tabular}{l|ccc|ccc|ccc}
    \toprule
    & \multicolumn{3}{c|}{Easy} & \multicolumn{3}{c|}{Medium} & \multicolumn{3}{c}{Hard} \\
    Model & PC1 Var & PC2 Var & $\rho$ & PC1 Var & PC2 Var & $\rho$ & PC1 Var & PC2 Var & $\rho$ \\
    \midrule
    SGNN + Objective & 1.000 & 0.000 & 0.600 & 1.000 & 0.000 & 0.000 & 1.000 & 0.000 & -0.200 \\
    MAGNN + Objective & 0.997 & 0.003 & 0.400 & 0.997 & 0.003 & -1.000 & 0.997 & 0.003 & -1.000 \\
    SGNN + Self-Predictor & 1.000 & 0.000 & 1.000 & 1.000 & 0.000 & -0.400 & 1.000 & 0.000 & -1.000 \\
    MAGNN + Self-Predictor & 0.990 & 0.010 & -0.200 & 0.989 & 0.011 & -1.000 & 0.989 & 0.011 & -1.000 \\
    \bottomrule
  \end{tabular}
  \caption{PCA diagnostics on final node embeddings. $\rho$ denotes Spearman correlation between original graph degree and PC1. Pearson correlations were not included in the current output and should be added in the next diagnostic run.}
  \label{tab:pca-diagnostics}
\end{table*}

The PCA diagnostics show that the embeddings are nearly one-dimensional in all variants, with PC1 explaining almost all variance. This is consistent with the existence of a dominant ranking direction. However, the sign and magnitude of Spearman correlation vary across regimes and models. Therefore, the diagnostic should not be interpreted as proof that PC1 is degree. Instead, it shows that the learned embeddings collapse onto a strong scalar ordering direction, which can be visualized in the placeholder figures above. The decisive result is not PC1 itself, but whether the learned updater can use the model's internal signals to support incremental residual decoding.

\section{Theoretical Motivation}

\subsection{Why Residual Decoding Matters}

Many algorithms for NP-hard and NP-complete problems do not make a single global decision. Instead, they proceed through a sequence of residual decisions: pruning elements, branching, decimating variables, traversing candidate states, or repeatedly simplifying the problem. This is not the only possible strategy, but it is a common and powerful algorithmic pattern.

For graph problems, a GNN pass on a dense graph costs $O(n^2)$. If one reruns the GNN after each of $O(n)$ residual decisions, the resulting decoder costs
\[
  O(n)\cdot O(n^2) = O(n^3).
\]
A natural way to reduce this cost is to run the GNN once and then perform $O(n)$ residual decisions, each with an $O(n)$ update. This gives
\[
  O(n^2) + O(n)\cdot O(n) = O(n^2).
\]
This approach combines the representation-learning ability of neural networks with the efficiency of classical iterative algorithms. The key question is whether the learned representation contains information that can be updated cheaply as the residual problem changes.

\subsection{Why Ranking Alone Is Not Enough}

\begin{lemma}[Ranking does not identify additive scale]
  Let $d_i$ be an additive statistic, such as residual degree. A training signal that only enforces the ordering of vertices identifies $d_i$ only up to a strictly monotone transformation.
\end{lemma}

\begin{proof}
  Suppose a score $s_i=d_i$ satisfies all ranking constraints. Let $\phi$ be any strictly increasing function and define $\tilde{s}_i=\phi(d_i)$. Then
  \[
    d_i>d_j \implies \phi(d_i)>\phi(d_j),
  \]
  so $\tilde{s}$ satisfies exactly the same ranking constraints. Therefore ranking supervision cannot distinguish $d_i$ from $\phi(d_i)$. In particular, it does not identify the additive scale required by residual updates.
\end{proof}

\begin{theorem}[Exact additive deletion updates require affine coordinates]
  Suppose a scalar representation $s_i=\phi(d_i)$ is updated after deleting a neighbor by a fixed additive rule. If the true statistic changes as $d_i\leftarrow d_i-1$, and the score update is exact for all $d_i$, then $\phi$ must be affine.
\end{theorem}

\begin{proof}
  Exact update compatibility requires a constant $c$ such that
  \[
    \phi(d-1)=\phi(d)-c
  \]
  for all relevant $d$. Therefore
  \[
    \phi(d)-\phi(d-1)=c
  \]
  for all $d$, so the discrete derivative of $\phi$ is constant. Hence $\phi(d)=ad+b$ for constants $a,b$. Thus only affine encodings preserve exact additive deletion updates.
\end{proof}

\subsection{Nonlinear Scores Create Residual Error}

\begin{theorem}[Non-affine coordinates induce update error]
  Let $s_i=\phi(d_i)$ for a non-affine function $\phi$. If one applies an additive residual update designed for $d_i$, then there exist values of $d$ for which the updated score differs from the correct residual score.
\end{theorem}

\begin{proof}
  The correct residual score after deleting one neighbor is $\phi(d-1)$. An additive update from $\phi(d)$ gives $\phi(d)-c$ for some constant $c$. If this were exact for all $d$, then by the previous theorem $\phi$ would have to be affine. Since $\phi$ is non-affine, there exists at least one $d$ such that
  \[
    \phi(d-1) \ne \phi(d)-c .
  \]
  Thus the update incurs error.
\end{proof}

This does not mean that nonlinear networks cannot rank vertices. Rather, it means that a nonlinear score that is good for ranking need not be a valid state variable for cheap residual updates.

\begin{theorem}[Local residual errors can propagate]
  Consider an iterative pruning procedure that removes the lowest-scoring vertex at each step. Suppose the estimated scores satisfy $\hat{s}_i=s_i+\epsilon_i$. If, for some pair $i,j$, the score gap satisfies
  \[
    |s_i-s_j| < |\epsilon_i-\epsilon_j|,
  \]
  then the estimated procedure may remove a different vertex than the exact procedure. After this mistake, the residual graphs differ, and all future scores are computed on different states.
\end{theorem}

\begin{proof}
  If $s_i<s_j$ but $\hat{s}_i>\hat{s}_j$, then the perturbed procedure reverses the removal order of $i$ and $j$. Once a different vertex is removed, the residual vertex set changes. Since subsequent scores depend on the residual graph, the future computation is no longer a small perturbation of the original trajectory; it is a trajectory on a different residual problem. Hence small local errors can cascade through the pruning process.
\end{proof}

\subsection{Why the Gradient Helps}

For the clique objective, the gradient is
\[
  \nabla_p \mathcal{L}_{\mathrm{clique}}(p;A,k)
  =
  -\frac{2Ap}{k(k-1)}
  +\frac{2Bp}{k(n-k)}
  +\frac{20(\sum_i p_i-k)}{n^2}\mathbf{1} .
\]
The terms $Ap$ and $Bp$ are linear graph marginals of the current relaxed solution. Thus, while the model is not given degree as an explicit semantic feature, the gradient exposes linear residual information derived from the objective.

\begin{proposition}[Gradient access makes additive marginals representable]
  For quadratic graph objectives of the form $p^\top A p$ and $p^\top Bp$, the gradient with respect to $p$ contains the linear marginals $Ap$ and $Bp$. Therefore a model that receives the objective gradient has direct access to the additive marginal quantities that change predictably under residual updates.
\end{proposition}

\begin{proof}
  For a symmetric matrix $A$,
  \[
    \nabla_p(p^\top A p) = 2Ap .
  \]
  Similarly, $\nabla_p(p^\top Bp)=2Bp$. These are linear functions of the current soft assignment $p$. Under a residual deletion, changing one coordinate of $p$ changes $Ap$ and $Bp$ by the corresponding adjacency and non-adjacency columns. Thus the gradient exposes exactly the type of local linear marginal signal needed by an incremental updater.
\end{proof}

The learned updater is therefore crucial. The gradient provides access to linear marginal information, but the updater learns how to transform the current score, removed vertex information, adjacency relation, and gradient channel into the next residual score. This avoids assuming that the GNN's output score itself is already degree-like.

\subsection{Main Theoretical Consequence}

The theory and experiments together support the following interpretation. A standard GNN can learn a useful ranking, and rerun-pruned decoding can exploit this ranking by recomputing the representation on every residual graph. However, this costs $O(n^3)$ on dense graphs. A cheap $O(n^2)$ incremental decoder requires more than a ranking: it requires update-relevant state information. The model-aware gradient channel exposes linear marginal information from the objective, and the learned updater uses this information to approximate rerun-style residual decoding without rerunning the full GNN at every step.
